\documentclass[../DoAn.tex]{subfiles}

\begin{document}

\section{Mục tiêu của luồng xử lý truy vấn}
\label{section:query-flow-goal}

Chương này trình bày luồng xử lý truy vấn trên dữ liệu đã được trích xuất từ PDF.
Nếu chương trước tập trung vào việc biến PDF thành block, chunk và index, thì
chương này tập trung vào cách hệ thống sử dụng các dữ liệu đó để tìm bằng chứng,
tạo câu trả lời và dẫn chứng nguồn. Vai trò của luồng truy vấn là kiểm chứng khả
năng khai thác dữ liệu đã trích xuất, không thay thế trọng tâm trích xuất PDF của
đề tài. Các kỹ thuật truy xuất và hỏi đáp được tích hợp trong luồng này dựa trên
các hướng BM25, dense retrieval, RRF và RAG đã trình bày ở Chương 3
\cite{robertson2009bm25,karpukhin2020dpr,cormack2009rrf,lewis2020rag}.

Đầu vào của luồng này là câu hỏi của người dùng và chỉ mục truy xuất đã được xây
từ kết quả ingest. Đầu ra là câu trả lời, danh sách bằng chứng, dẫn chứng và vết xử
lý. Các thành phần truy xuất và hỏi đáp được tách riêng để có thể đánh giá lỗi
theo từng tầng.

\section{Tổng quan luồng xử lý truy vấn}
\label{section:query-flow-overview}

\begin{figure}[H]
\centering
\small
\begin{tabular}{c}
\fbox{\parbox{0.82\linewidth}{\centering Câu hỏi của người dùng}}\\[0.15cm]
\(\downarrow\)\\[-0.05cm]
\fbox{\parbox{0.82\linewidth}{\centering Bộ định tuyến câu hỏi: xác định loại câu hỏi và tín hiệu bảng nếu có}}\\[0.15cm]
\(\downarrow\)\\[-0.05cm]
\fbox{\parbox{0.82\linewidth}{\centering Bộ lập kế hoạch truy xuất: chọn chiến lược, top-k, candidate-k, sắp xếp lại}}\\[0.15cm]
\(\downarrow\)\\[-0.05cm]
\fbox{\parbox{0.82\linewidth}{\centering Bộ truy xuất: BM25, dense, hybrid; điều chỉnh theo metadata bảng khi được bật}}\\[0.15cm]
\(\downarrow\)\\[-0.05cm]
\fbox{\parbox{0.82\linewidth}{\centering Bộ kiểm tra bằng chứng: kiểm tra liên quan, đầy đủ và khả năng dẫn chứng}}\\[0.15cm]
\(\downarrow\)\\[-0.05cm]
\fbox{\parbox{0.82\linewidth}{\centering Bộ sinh câu trả lời: tạo câu trả lời có căn cứ và dẫn chứng}}\\[0.15cm]
\(\downarrow\)\\[-0.05cm]
\fbox{\parbox{0.82\linewidth}{\centering Kết quả hỏi đáp: câu trả lời, bằng chứng, dẫn chứng, vết xử lý, thời gian}}\\
\end{tabular}
\caption{Luồng xử lý truy vấn trên dữ liệu PDF đã trích xuất}
\label{fig:query-flow-overview}
\end{figure}

Luồng này được thiết kế để câu trả lời cuối cùng luôn gắn với bằng chứng trong
chỉ mục. Với câu hỏi bảng, hệ thống có thể ưu tiên các chunk ở mức hàng hoặc ô.
Với câu hỏi văn bản, hệ thống ưu tiên các đoạn có nội dung và ngữ cảnh đề mục
phù hợp.

\section{Liên kết yêu cầu với thiết kế truy vấn}
\label{section:query-traceability}

Bảng~\ref{tab:query-traceability} tóm tắt các yêu cầu liên quan đến khai thác dữ
liệu PDF sau trích xuất. Các yêu cầu này sử dụng trực tiếp corpus và metadata do
luồng ingest tạo ra.

\begin{table}[H]
\centering
\small
\renewcommand{\arraystretch}{1.2}
\caption{Liên kết yêu cầu với thành phần thiết kế trong luồng truy vấn}
\label{tab:query-traceability}
\begin{tabularx}{\textwidth}{
    |>{\raggedright\arraybackslash}p{0.16\textwidth}
    |>{\raggedright\arraybackslash}X
    |>{\raggedright\arraybackslash}p{0.30\textwidth}|
}
\hline
\textbf{Mã} & \textbf{Yêu cầu liên quan} & \textbf{Thành phần thiết kế} \\
\hline
FR-05 &
Truy xuất các đoạn thông tin liên quan đến câu hỏi của người dùng. &
\texttt{RetrievalService}, BM25, dense retrieval, hybrid retrieval. \\
\hline
FR-05, FR-07 &
Ưu tiên truy xuất và dẫn chứng chính xác hơn với câu hỏi liên quan đến bảng. &
\texttt{QueryRouter}, table-aware retrieval, metadata bảng/ô. \\
\hline
FR-06 &
Tạo câu trả lời dựa trên bằng chứng đã truy xuất và tránh tự suy đoán khi thiếu
thông tin. &
\texttt{EvidenceChecker}, \texttt{GroundedAnswerGenerator}. \\
\hline
FR-07 &
Cung cấp dẫn chứng về tài liệu, trang, đoạn văn, bảng, hàng hoặc ô bảng. &
\texttt{citations.py}, metadata dẫn chứng. \\
\hline
FR-08 &
Hỗ trợ đánh giá thực nghiệm theo từng tầng. &
Vết xử lý, bằng chứng được chọn, chiến lược truy xuất, thời gian xử lý, kết quả benchmark. \\
\hline
\end{tabularx}
\end{table}

\section{Dữ liệu vào và đầu ra của luồng truy vấn}
\label{section:query-input-output}

\begin{table}[H]
\centering
\small
\renewcommand{\arraystretch}{1.2}
\caption{Dữ liệu vào và đầu ra của luồng xử lý truy vấn}
\label{tab:query-input-output}
\begin{tabularx}{\textwidth}{
    |>{\raggedright\arraybackslash}p{0.26\textwidth}
    |>{\raggedright\arraybackslash}X
    |>{\raggedright\arraybackslash}p{0.28\textwidth}|
}
\hline
\textbf{Thành phần} & \textbf{Dữ liệu chính} & \textbf{Vai trò} \\
\hline
Câu hỏi &
Câu hỏi ngôn ngữ tự nhiên, có thể là câu hỏi văn bản hoặc tra cứu bảng. &
Kích hoạt luồng truy xuất và hỏi đáp. \\
\hline
Chỉ mục truy xuất &
Corpus chunk, chỉ mục BM25, embedding nếu có, metadata nguồn. &
Không gian tìm kiếm bằng chứng. \\
\hline
Cấu hình truy xuất &
Chiến lược truy xuất, top-k, candidate-k, trọng số BM25/dense, bộ sắp xếp lại. &
Điều khiển cách lấy bằng chứng. \\
\hline
RetrievedHit &
Chunk được truy xuất, điểm, thứ hạng, nội dung và metadata nguồn. &
Đơn vị bằng chứng đầu vào cho QA. \\
\hline
Đánh giá bằng chứng &
Điểm liên quan, độ bao phủ, mức có căn cứ và trạng thái đánh giá bằng chứng. &
Quyết định sinh câu trả lời hoặc trả lời thận trọng khi chưa đủ căn cứ. \\
\hline
QAResult &
Câu trả lời, dẫn chứng, bằng chứng được chọn, chiến lược, vết xử lý và thời gian xử lý. &
Đầu ra cuối cùng cho người dùng và benchmark. \\
\hline
\end{tabularx}
\end{table}

\section{Phân loại câu hỏi và lập kế hoạch truy xuất}
\label{section:query-router-planner-design}

Hệ thống sử dụng QueryRouter để xác định loại câu hỏi ở mức đủ phục vụ retrieval.
Bộ phân loại hiện tại dựa trên luật và từ khóa, không phải mô hình học máy. Cách
làm này phù hợp với phạm vi đồ án vì dễ kiểm tra, dễ giải thích và dễ tái lập trong
benchmark. Khác với các hướng adaptive RAG học chiến lược truy xuất bằng mô
hình \cite{jeong2024adaptiverag}, thiết kế hiện tại ưu tiên tính kiểm soát và khả
năng phân tích lỗi.

Sau khi có loại câu hỏi, \texttt{QueryAwareRetrievalPlanner} chọn cấu hình truy
xuất. Ví dụ, câu hỏi dạng tra cứu bảng mở rộng tập ứng viên và tăng trọng số
BM25. Khi tùy chọn table-aware retrieval được bật, bộ sắp xếp lại mới cộng thêm
điểm dựa trên metadata bảng, hàng và ô. Câu hỏi cần tổng hợp có thể bật bước sắp
xếp lại kết quả, còn câu hỏi đơn giản dùng top-k nhỏ hơn để giảm nhiễu.

\begin{table}[H]
\centering
\small
\renewcommand{\arraystretch}{1.2}
\caption{Thiết kế truy xuất theo loại câu hỏi}
\label{tab:query-type-retrieval-plan}
\begin{tabularx}{\textwidth}{
    |>{\raggedright\arraybackslash}p{0.24\textwidth}
    |>{\raggedright\arraybackslash}p{0.26\textwidth}
    |>{\raggedright\arraybackslash}X|
}
\hline
\textbf{Loại câu hỏi} & \textbf{Cấu hình ưu tiên} & \textbf{Lý do thiết kế} \\
\hline
Factoid / definition &
Hybrid, top-k vừa phải. &
Cần tìm đoạn chứa thông tin trực tiếp. \\
\hline
Policy / procedural &
Hybrid + sắp xếp lại. &
Cần giữ nhiều evidence vì thường có điều kiện hoặc quy trình. \\
\hline
Comparison / multi-hop &
Hybrid hoặc RRF, candidate-k lớn hơn. &
Cần thu thập evidence từ nhiều đoạn. \\
\hline
Table lookup &
Hybrid, tăng trọng số BM25; có thể bật table-aware boost khi cấu hình cho phép. &
Cần khớp chính xác hàng, cột, giá trị số hoặc ký hiệu trong bảng. \\
\hline
Ambiguous &
Hybrid với phạm vi tìm kiếm rộng hơn. &
Ý định chưa rõ nên không thu hẹp quá sớm. \\
\hline
\end{tabularx}
\end{table}

\section{Truy xuất evidence}
\label{section:evidence-retrieval-design}

\texttt{RetrievalService} cung cấp một giao diện chung cho nhiều chiến lược truy xuất.
BM25 khai thác tín hiệu từ khóa \cite{robertson2009bm25}; dense retrieval khai
thác embedding \cite{karpukhin2020dpr,reimers2019sentencebert}; hybrid
retrieval kết hợp hai tín hiệu này bằng weighted sum hoặc RRF
\cite{cormack2009rrf}. Reranker chỉ áp dụng trên tập ứng viên nhỏ đã được lấy ra
trước đó. Trong các cấu hình sắp xếp lại nâng cao, ColBERT là một ví dụ của
hướng late interaction giữa truy vấn và tài liệu \cite{khattab2020colbert}.

\begin{figure}[H]
\centering
\small
\begin{tabular}{c}
\fbox{\parbox{0.82\linewidth}{\centering Câu hỏi + cấu hình truy xuất}}\\[0.15cm]
\(\downarrow\)\\[-0.05cm]
\fbox{\parbox{0.82\linewidth}{\centering Truy xuất ứng viên: BM25, dense hoặc hybrid}}\\[0.15cm]
\(\downarrow\)\\[-0.05cm]
\fbox{\parbox{0.82\linewidth}{\centering Sắp xếp lại tùy chọn: heuristic, cross-encoder hoặc ColBERT}}\\[0.15cm]
\(\downarrow\)\\[-0.05cm]
\fbox{\parbox{0.82\linewidth}{\centering Điều chỉnh theo metadata bảng nếu truy vấn liên quan bảng}}\\[0.15cm]
\(\downarrow\)\\[-0.05cm]
\fbox{\parbox{0.82\linewidth}{\centering Top-k kết quả có nội dung, điểm và metadata nguồn}}\\
\end{tabular}
\caption{Luồng truy xuất evidence trong hệ thống}
\label{fig:evidence-retrieval-flow}
\end{figure}

Mỗi kết quả truy xuất được chuẩn hóa thành \texttt{RetrievedHit}. Đối tượng này
không chỉ chứa nội dung chunk, mà còn giữ điểm truy xuất, thứ hạng và metadata
như tài liệu nguồn, trang, block type, heading path và thông tin bảng nếu có. Đây
là dữ liệu mà tầng QA sử dụng để kiểm tra evidence và tạo citation.

\section{Truy xuất trên dữ liệu bảng}
\label{section:query-table-aware-retrieval}

Truy vấn bảng là nơi thể hiện rõ vai trò của dữ liệu đã được trích xuất ở chương
trước. Nếu bảng chỉ được đưa vào index như văn bản phẳng, hệ thống có thể tìm
đúng bảng nhưng khó xác định đúng hàng hoặc ô. Vì vậy, khi câu hỏi có dấu hiệu
tra cứu bảng, hệ thống ưu tiên các chunk có metadata bảng.
Thiết kế này kế thừa phần biểu diễn bảng ở luồng ingest, trong đó cấu trúc bảng
có thể được tái dựng từ Table Transformer/TATR và các hộp từ lấy từ lớp văn bản
PDF~\cite{smock2022pubtables}. Với tài liệu scan, OCR là nguồn văn bản và tọa
độ bổ sung trong hướng tích hợp riêng~\cite{du2020ppocr}.

\begin{table}[H]
\centering
\small
\renewcommand{\arraystretch}{1.2}
\caption{Metadata bảng được dùng trong truy xuất và dẫn chứng}
\label{tab:table-query-metadata}
\begin{tabularx}{\textwidth}{
    |>{\raggedright\arraybackslash}p{0.30\textwidth}
    |>{\raggedright\arraybackslash}X|
}
\hline
\textbf{Trường} & \textbf{Vai trò} \\
\hline
\texttt{table\_id} &
Xác định bảng nguồn của evidence. \\
\hline
\texttt{row\_header} &
Hỗ trợ khớp câu hỏi với hàng tương ứng trong bảng. \\
\hline
\texttt{col\_header} &
Hỗ trợ khớp câu hỏi với cột hoặc điều kiện cần tra cứu. \\
\hline
\texttt{cell\_text} &
Giá trị có thể được dùng trực tiếp cho câu trả lời bảng. \\
\hline
\texttt{citation\_target} &
Xác định citation ở mức bảng, hàng hoặc ô. \\
\hline
\end{tabularx}
\end{table}

Cách xử lý này giúp hệ thống trả về evidence nhỏ hơn và chính xác hơn. Ví dụ,
với câu hỏi về một giá trị trong bảng điểm, evidence phù hợp nhất thường là
table cell chứa \texttt{row\_header}, \texttt{col\_header} và \texttt{cell\_text}, không
phải toàn bộ trang PDF.

\section{Kiểm tra evidence và sinh câu trả lời}
\label{section:evidence-check-answer-design}

Sau truy xuất, EvidenceChecker đánh giá tập evidence theo mức độ liên quan, độ
đầy đủ, khả năng hỗ trợ câu trả lời và khả năng tạo citation. Kết quả kiểm tra được
biểu diễn bằng bốn trạng thái: \texttt{answer}, \texttt{expand\_retrieval},
\texttt{switch\_strategy} và \texttt{abstain}. Trong triển khai hiện tại, các trạng
thái này được lưu trong vết xử lý và dùng để kiểm soát bộ sinh câu trả lời. Hệ
thống chưa tự động thực hiện lại truy xuất khi nhận trạng thái
\texttt{expand\_retrieval} hoặc \texttt{switch\_strategy}. Cách tổ chức này đi
theo hướng RAG có kiểm soát bằng chứng~\cite{lewis2020rag,asai2024selfrag}.

\begin{table}[H]
\centering
\small
\renewcommand{\arraystretch}{1.2}
\caption{Các trạng thái sau bước kiểm tra evidence}
\label{tab:evidence-check-decisions-revised}
\begin{tabularx}{\textwidth}{
    |>{\raggedright\arraybackslash}p{0.24\textwidth}
    |>{\raggedright\arraybackslash}X|
}
\hline
\textbf{Quyết định} & \textbf{Xử lý trong hệ thống} \\
\hline
\texttt{answer} &
Evidence đủ để sinh câu trả lời có căn cứ. \\
\hline
\texttt{expand\_retrieval} &
Evidence có liên quan nhưng chưa đủ bao phủ câu hỏi; hệ thống ghi nhận trạng
thái và không tạo câu trả lời khẳng định. \\
\hline
\texttt{switch\_strategy} &
Evidence truy xuất còn yếu; hệ thống ghi nhận trạng thái và không tự động đổi
chiến lược trong luồng hiện tại. \\
\hline
\texttt{abstain} &
Không đủ căn cứ; hệ thống trả lời thận trọng thay vì suy đoán. \\
\hline
\end{tabularx}
\end{table}

AnswerGenerator tạo câu trả lời từ các evidence đã chọn. Với câu hỏi bảng, hệ
thống có thể ưu tiên giá trị trong \texttt{cell\_text}. Với câu hỏi văn bản, hệ thống
tạo câu trả lời dựa trên các câu hoặc đoạn được evidence hỗ trợ. Các thành phần
LLM như fallback hoặc explainer chỉ là tùy chọn: chúng không được dùng để thay
thế retrieval và evidence checking \cite{lewis2020rag}.

\section{Thiết kế citation và trace}
\label{section:citation-trace-design}

Citation được tạo từ metadata của evidence. Với văn bản, citation trỏ về tài liệu,
trang, chunk và heading path. Với bảng, citation có thể trỏ về table id, hàng, cột,
ô và bbox nguồn. Trace lưu lại loại câu hỏi, chiến lược truy xuất, top hit, quyết
định evidence và nguồn sinh câu trả lời. Nhờ đó, khi kết quả sai, có thể kiểm tra
lỗi đến từ retrieval, evidence checking hay answer generation.

\begin{table}[H]
\centering
\small
\renewcommand{\arraystretch}{1.2}
\caption{Các mức dẫn chứng được hỗ trợ}
\label{tab:citation-levels-query}
\begin{tabularx}{\textwidth}{
    |>{\raggedright\arraybackslash}p{0.24\textwidth}
    |>{\raggedright\arraybackslash}X|
}
\hline
\textbf{Mức citation} & \textbf{Thông tin dẫn chứng} \\
\hline
Text chunk &
Tài liệu, trang, heading path, chunk id. \\
\hline
Table &
Trang, table id, caption hoặc mô tả bảng. \\
\hline
Table row &
Table id, chỉ số hàng, row header và trang. \\
\hline
Table cell &
Table id, row header, column header, cell text và bbox nếu có. \\
\hline
\end{tabularx}
\end{table}

\section{Tổ chức mã nguồn của luồng truy vấn}
\label{section:query-code-organization}

\begin{table}[H]
\centering
\small
\renewcommand{\arraystretch}{1.2}
\caption{Các thành phần mã nguồn chính của luồng xử lý truy vấn}
\label{tab:query-code-organization}
\begin{tabularx}{\textwidth}{
    |>{\raggedright\arraybackslash}p{0.36\textwidth}
    |>{\raggedright\arraybackslash}X|
}
\hline
\textbf{Thành phần} & \textbf{Vai trò} \\
\hline
\texttt{app/retrieval} &
BM25, dense retrieval, hybrid retrieval, reranking, route planner và service. \\
\hline
\texttt{app/retrieval/table\_aware.py} &
Điều chỉnh truy xuất với các chunk thuộc bảng, hàng hoặc ô. \\
\hline
\texttt{app/qa/router.py} &
Phân loại câu hỏi bằng luật để chọn hướng xử lý phù hợp. \\
\hline
\texttt{app/qa/evidence\_checker.py} &
Đánh giá evidence, tạo trạng thái kiểm tra và cung cấp dữ liệu cho bộ sinh câu
trả lời. \\
\hline
\texttt{app/qa/answer\_generator.py} &
Sinh câu trả lời từ evidence đã chọn và ưu tiên câu trả lời có căn cứ. \\
\hline
\texttt{app/qa/citations.py} &
Chuẩn hóa citation từ metadata của evidence. \\
\hline
\texttt{scripts/quick\_qa\_terminal.py} &
Demo hỏi đáp nhanh trên terminal và hiển thị evidence/citation. \\
\hline
\end{tabularx}
\end{table}

\section{Tổng kết chương}
\label{section:query-chapter-summary}

Chương này đã trình bày luồng xử lý truy vấn trên dữ liệu đã trích xuất từ PDF.
Luồng này gồm phân loại câu hỏi, lập kế hoạch truy xuất, lấy evidence, kiểm tra
evidence, sinh câu trả lời và tạo citation. Cấu trúc này giúp hệ thống không chỉ
trả lời câu hỏi, mà còn chỉ ra dữ liệu PDF nào đã hỗ trợ câu trả lời đó. Chương
thực nghiệm tiếp theo sẽ đánh giá riêng từng tầng để chứng minh vai trò của ingest,
retrieval, table-aware chunking và grounded QA.

\end{document}
